% Generated draft: requires mathematical and source review.
\documentclass[11pt]{article}
\usepackage[margin=25mm]{geometry}
\usepackage{fontspec}
\setmainfont{Latin Modern Roman}
\usepackage{amsmath,amssymb,mathtools,unicode-math}
\setmathfont{Latin Modern Math}
\usepackage{enumitem,needspace,xurl,hyperref}
\setlength{\parindent}{0pt}
\setlength{\parskip}{5pt}
\setlength{\emergencystretch}{3em}
\newcommand{\N}{\mathbb{N}}
\newcommand{\Z}{\mathbb{Z}}
\newcommand{\Q}{\mathbb{Q}}
\newcommand{\R}{\mathbb{R}}
\newcommand{\C}{\mathbb{C}}
\newcommand{\F}{\mathbb{F}}
\newcommand{\A}{\mathbb{A}}
\newcommand{\PP}{\mathbb{P}}
\newcommand{\E}{\mathbb{E}}
\newcommand{\eps}{\varepsilon}
\newcommand{\dd}{\,\mathrm{d}}
\begin{document}
\section*{UnsolvedMath ID: 30004796}
\textbf{Source problem code:} catalogue entry
\subsection{Definitions and mathematical statement}
The supplied record calls the following assertion the refined formal-degree conjecture. Write \(A_G\) for the split center of the group \(G\), and \(\deg(\pi,\mu_G)\) for the formal degree of the representation \(\pi\) with respect to the measure \(\mu_G\). The source also uses a parameter \(\phi\), a representation \(\rho\), a group \(S_\phi\), an invariant \(r(G)\), an adjoint representation \(\operatorname{Ad}_G\), and a principal parameter \(\phi_{\operatorname{prin}}\). Their precise definitions and admissibility conditions are not supplied; these symbols retain the intended source meanings without imposing additional hypotheses.

For the intended data, the proposed identity under the hypothesis \(A_G=\{1\}\) is
\[
(-1)^{r(G)}\deg(\pi,\mu_G)
=
\sigma\,\frac{\dim\rho}{|S_\phi|}
\frac{\gamma\!\left(0,\operatorname{Ad}_G\circ\phi\right)}
{\gamma\!\left(0,\operatorname{Ad}_G\circ\phi_{\operatorname{prin}}\right)},
\qquad \sigma\in\{+1,-1\}.
\]
Here \(\gamma\) denotes the gamma factor appearing in the supplied formula. The source says that \(\sigma\) is given by the quotient of the corresponding root numbers. It does not explicitly define that quotient or the normalizations of the root numbers and gamma factors.

This is a conditional transcription requiring further specification, rather than a fully quantified conjecture. In particular, the source does not specify the base field, the class of groups or representations, the relation between \(\pi\) and \((\phi,\rho)\), the meaning of \(S_\phi\) and \(r(G)\), the normalization of \(\mu_G\), or the interpretation of the displayed quotient at \(s=0\) if its factors are singular or vanish. The intended range of quantification and these conventions must be recovered before a precise mathematical formulation can be accepted.
\subsection{Short English statement}
When the split center is trivial, the conjecture relates a signed formal degree to a representation-dimension factor and a ratio of adjoint gamma factors, with a sign prescribed by root numbers. Essential hypotheses and normalizations are missing from the supplied excerpt. The supplied dated curation labels the problem partially solved and distinguishes a claimed theorem for tempered orthogonal root numbers from the partly resolved refined formal-degree statement; these assessments are not independently verified here.
\subsection{Sources}
\begin{itemize}
\item Supplied local UnsolvedMath registry record: permanent ID \(30004796\), source problem code \texttt{catalogue entry}, title \emph{The Refined Formal-Degree Conjecture}, in the Oberwolfach Reports Open Problems collection. No individual proposer is supplied; the proposed year is recorded as 2021.
\item Supplied source citation: \emph{Automorphic Forms, Geometry and Arithmetic (2022)}, Oberwolfach Reports, pp.\ 2089--2156, DOI \texttt{10.4171/owr/2021/39}. Original Oberwolfach report: \url{https://doi.org/10.4171/owr/2021/39}. The background describes these reports as documenting mathematical workshops and their open-problem discussions.
\item Supplied literature reference: Schwein, orthogonal root numbers, \url{https://arxiv.org/abs/2107.02360}.
\item Supplied literature reference, labeled ``published orthogonal-root-number theorem'': \url{https://doi.org/10.1007/s00208-022-02416-6}.
\item Literature assessment dated exactly \texttt{2026-08-22}: the supplied curation states that the refined central-character identity for orthogonal root numbers is proved for tempered parameters, while the accompanying refined formal-degree statement and its full generality remain only partly resolved. This is supplied dated curation derived from the OWR record, not an independently verified literature assessment. Its time-sensitive status and cited claims require independent verification.
\item The registry labels the statement \texttt{corrected\_verified}, with the explanation that explicit formulas were dollar-delimited and the supplied mathematical content was preserved. This records an editorial check and does not establish mathematical correctness or completeness.
\end{itemize}
\end{document}
